\chapter{Conclusiones}
\label{cap:conclusiones-final}

\section{Principales logros}

\begin{enumerate}[nosep]
  \item Un sistema web funcional de gesti\'on veterinaria, con
    autenticaci\'on y autorizaci\'on por rol/propiedad, CRUD de mascotas, y
    los m\'odulos de citas, consultas y vacunas de la Unidad~IV,
    respaldados por seis rutinas almacenadas de PostgreSQL con acceso JPA
    formal.
  \item Evidencia emp\'irica multim\'etrica real: 205 pruebas
    automatizadas (0 fallos), cobertura JaCoCo 91{,}80\,\%/79{,}39\,\% sobre
    un umbral autom\'atico del 70\,\%, diez corridas de rendimiento con
    0{,}0\,\% de tasa de error, usabilidad SUS medida con 18 participantes
    reales, y auditor\'ia de seguridad por tres m\'etodos independientes.
  \item Despliegue real en producci\'on (Render, HTTPS activo,
    \emph{healthcheck} \texttt{UP}), imagen de contenedor inmutable
    publicada en GHCR, y archivado permanente en Zenodo con DOI reales
    para el software y el conjunto de datos de evidencias.
  \item Cierre completo (15/15) de las observaciones docentes formales
    acumuladas a lo largo de tres entregas previas.
\end{enumerate}

\section{Evidencia}

Cada afirmaci\'on cuantitativa de este informe cita su fuente real en el
repositorio (archivo, script o reporte espec\'ifico), reproducible por un
tercero con los comandos documentados en el
cap\'itulo~\ref{cap:despliegue}. Ninguna cifra se estim\'o ni se
extrapol\'o de una entrega anterior sin verificaci\'on directa contra la
evidencia vigente.

\section{Valor del proyecto}

M\'as all\'a de la funcionalidad implementada, el valor distintivo de
BIOPET, frente al panorama de trabajos relacionados revisado
(cap\'itulo~\ref{cap:trabajos-relacionados}), es la \textbf{amplitud
verificable de su evidencia emp\'irica} sobre un mismo artefacto: ning\'un
trabajo comparado reporta simult\'aneamente cobertura de pruebas,
rendimiento, usabilidad, calidad web y seguridad por tres \'angulos
distintos, todo enmarcado bajo ISO/IEC~25010 y con reproducibilidad de
infraestructura de punta a punta. Esto no implica superioridad sobre esos
trabajos: implica que la pr\'actica de documentaci\'on y evidencia aplicada
en BIOPET permite razonar sobre relaciones entre dimensiones de calidad
(por ejemplo, entre cobertura de pruebas y ausencia de hallazgos de
seguridad est\'atica) que un estudio de una sola dimensi\'on no puede
mostrar.

\section{Estado de release}

BIOPET \texttt{v1.0.0} representa el cierre funcional y de evidencia de la
Entrega Final, pero \textbf{no debe interpretarse como un sistema listo
para producci\'on industrial sin condiciones}: el despliegue actual usa el
plan gratuito de Render (con purgado de datos a los 30 d\'ias, sin
verificaci\'on posible hoy de operaci\'on sostenida m\'as all\'a de ese
horizonte), el rate limiting de login no es distribuido, no existe
integraci\'on con un SIEM, y varios requisitos \emph{Should}/\emph{Could}
heredados de la Entrega~1A permanecen pendientes. El tag Git \texttt{v1.0.0}
y el GitHub Release correspondiente todav\'ia no se han creado a la fecha
de cierre de este informe (fuera del alcance expl\'icito de esta tarea de
redacci\'on).

\section{Trabajo futuro}

\begin{itemize}[nosep]
  \item Completar el historial cl\'inico longitudinal, el calendario
    interactivo de citas, la facturaci\'on digital y los reportes
    exportables (\texttt{REQ-F-013}, \texttt{REQ-F-015}, \texttt{REQ-F-018},
    \texttt{REQ-F-019}).
  \item Redactar la narrativa completa del \texttt{README.md} de
    \texttt{docs/mediciones/lighthouse/} para la corrida mobile+desktop del
    2026-08-18 (ya archivada con JSON crudos y las cuatro categor\'ias
    cumplidas), documentando expl\'icitamente el mecanismo de sesi\'on que
    permiti\'o auditar \texttt{/mascotas} directamente en esa corrida.
  \item Medir rendimiento bajo carga espec\'ificamente para los endpoints
    de citas, consultas y vacunas, no solo para el listado de mascotas.
  \item Migrar a un plan de base de datos persistente (o instrumentar
    verificaci\'on de respaldo/restauraci\'on) para poder evaluar
    emp\'iricamente la operaci\'on sostenida m\'as all\'a de 30 d\'ias.
  \item Evaluar un mecanismo de rate limiting distribuido si el sistema
    llegara a desplegarse con m\'ultiples r\'eplicas del backend.
  \item Integrar los logs \texttt{AUTH\_AUDIT} con un sistema de
    recolecci\'on y alertamiento centralizado (SIEM).
  \item Documentar el conteo de candidatos evaluados y excluidos del
    bloque de Jaime en la revisi\'on PRISMA, para homogeneizar el nivel de
    detalle entre los tres bloques (cap\'itulo~\ref{cap:trabajos-relacionados}).
  \item Crear el tag Git \texttt{v1.0.0} y el GitHub Release
    correspondiente, lo que activar\'a autom\'aticamente la publicaci\'on
    de las etiquetas \texttt{1.0.0}/\texttt{latest} en GHCR.
\end{itemize}
